\section{Generating the Single-Qutrit Clifford Group}
\label{sec:clifford-gen}

Both the physical RB baseline and the LRB-D protocol draw every Clifford uniformly from the single-qutrit Clifford group. We therefore need an explicit catalog of its elements, a realization of each element as a circuit, and the inverse of each element for the mirror inverse $C_m^{-1}, \dots, C_1^{-1}$, applied in this order. This appendix records how that catalog is constructed and how its elements become circuits.

Modulo global phases, the single-qutrit Clifford group $\mathcal{C}_3$ contains $216$ elements. It factors into the qutrit Pauli group and the quotient $\mathcal{C}_3 / \mathcal{P}_3$, which is isomorphic to the symplectic group $Sp(2,\mathbb{Z}_3) = SL(2,\mathbb{Z}_3)$ of order $24$ \cite{Appleby2005,Hostens2005}. Counting phase-free Pauli representatives $X^a Z^b$ with $a,b \in \mathbb{Z}_3$ gives $9$ elements, and $9 \times 24 = 216$. We exploit this factorization and build the group in two stages, first the $24$ coset representatives and then their products with the $9$ Pauli monomials.

\paragraph{Coset representatives from a Cayley graph.}
The quotient $\mathcal{C}_3/\mathcal{P}_3$ is generated by the images of $H$ and $S$. Two Cliffords lie in the same coset exactly when they act in the same way, by conjugation, on the eight non-identity Pauli classes $\{X, Z, XZ, X^2, Z^2, X^2Z, Z^2X, X^2Z^2\}$. This action therefore realizes the quotient faithfully as a permutation group on eight symbols. We perform a breadth-first traversal of the Cayley graph of this group. Starting from the identity, every element of the current level is multiplied on the right by $H$ and then by $S$, and the first time a new element is produced we record its label as the label of its parent extended by the corresponding generator symbol. Because the traversal is breadth-first, each of the $24$ elements receives a minimal-length word in $H$ and $S$, and the traversal terminates once all $24$ are found. The longest of these words has five symbols. In the figure and in the text we write repeated symbols as powers, so for instance $SSHS$ appears as $S^2HS$. The power notation is only a label, and in the circuits every symbol remains a separate gate. The graph explored by this traversal is shown in \autoref{fig:cayley-c3modp3}, with every node carrying the word that the traversal assigned to it.

\begin{figure}[H]
  \centering
  \resizebox{0.92\textwidth}{!}{%
\begin{tikzpicture}[>={Stealth[length=2.2mm,width=1.6mm]},
  cnode/.style={circle,draw=black,fill=white,inner sep=1.1pt,minimum size=7.5mm,font=\scriptsize},
  Hedge/.style={->,HCol,line width=0.75pt},
  Sedge/.style={->,SwapCol,line width=0.75pt}]
  \node[cnode] (n0) at (2.83,5.09) {$I$};
  \node[cnode] (n1) at (-1.32,6.5) {$H$};
  \node[cnode] (n2) at (2.41,0.84) {$S$};
  \node[cnode] (n3) at (1.74,6.4) {$H^2$};
  \node[cnode] (n4) at (-0.84,2.41) {$HS$};
  \node[cnode] (n5) at (-2.41,-0.84) {$SH$};
  \node[cnode] (n6) at (6.5,1.32) {$S^2$};
  \node[cnode] (n7) at (-2.49,5.26) {$H^3$};
  \node[cnode] (n8) at (1.0,2.35) {$H^2S$};
  \node[cnode] (n9) at (2.35,-1.0) {$HSH$};
  \node[cnode] (n10) at (-5.09,2.83) {$HS^2$};
  \node[cnode] (n11) at (-6.5,-1.32) {$SHS$};
  \node[cnode] (n12) at (5.09,-2.83) {$S^2H$};
  \node[cnode] (n13) at (-2.35,1.0) {$H^3S$};
  \node[cnode] (n14) at (-1.0,-2.35) {$H^2SH$};
  \node[cnode] (n15) at (5.26,2.49) {$H^2S^2$};
  \node[cnode] (n16) at (6.4,-1.74) {$HSHS$};
  \node[cnode] (n17) at (-6.4,1.74) {$SHSH$};
  \node[cnode] (n18) at (-2.83,-5.09) {$SHS^2$};
  \node[cnode] (n19) at (1.32,-6.5) {$S^2HS$};
  \node[cnode] (n20) at (0.84,-2.41) {$H^3SH$};
  \node[cnode] (n21) at (-5.26,-2.49) {$H^2SHS$};
  \node[cnode] (n22) at (2.49,-5.26) {$HSHS^2$};
  \node[cnode] (n23) at (-1.74,-6.4) {$S^2HSH$};
  \draw[Hedge] (n0) to[bend left=16] (n1);
  \draw[Sedge] (n0) to[bend left=16] (n2);
  \draw[Hedge] (n1) to[bend left=10] (n3);
  \draw[Sedge] (n1) to[bend left=16] (n4);
  \draw[Hedge] (n2) to[bend left=16] (n5);
  \draw[Sedge] (n2) to[bend left=16] (n6);
  \draw[Hedge] (n3) to[bend left=16] (n7);
  \draw[Sedge] (n3) to[bend left=16] (n8);
  \draw[Hedge] (n4) to[bend left=16] (n9);
  \draw[Sedge] (n4) to[bend left=16] (n10);
  \draw[Hedge] (n5) to[bend left=16] (n8);
  \draw[Sedge] (n5) to[bend left=16] (n11);
  \draw[Hedge] (n6) to[bend left=16] (n12);
  \draw[Sedge] (n6) to[bend left=16] (n0);
  \draw[Hedge] (n7) to[bend left=16] (n0);
  \draw[Sedge] (n7) to[bend left=16] (n13);
  \draw[Hedge] (n8) to[bend left=16] (n14);
  \draw[Sedge] (n8) to[bend left=16] (n15);
  \draw[Hedge] (n9) to[bend left=16] (n13);
  \draw[Sedge] (n9) to[bend left=16] (n16);
  \draw[Hedge] (n10) to[bend left=16] (n11);
  \draw[Sedge] (n10) to[bend left=16] (n1);
  \draw[Hedge] (n11) to[bend left=10] (n17);
  \draw[Sedge] (n11) to[bend left=16] (n18);
  \draw[Hedge] (n12) to[bend left=16] (n15);
  \draw[Sedge] (n12) to[bend left=16] (n19);
  \draw[Hedge] (n13) to[bend left=16] (n20);
  \draw[Sedge] (n13) to[bend left=16] (n17);
  \draw[Hedge] (n14) to[bend left=16] (n2);
  \draw[Sedge] (n14) to[bend left=16] (n21);
  \draw[Hedge] (n15) to[bend left=16] (n16);
  \draw[Sedge] (n15) to[bend left=16] (n3);
  \draw[Hedge] (n16) to[bend left=10] (n6);
  \draw[Sedge] (n16) to[bend left=16] (n22);
  \draw[Hedge] (n17) to[bend left=16] (n21);
  \draw[Sedge] (n17) to[bend left=16] (n7);
  \draw[Hedge] (n18) to[bend left=16] (n19);
  \draw[Sedge] (n18) to[bend left=16] (n5);
  \draw[Hedge] (n19) to[bend left=10] (n23);
  \draw[Sedge] (n19) to[bend left=16] (n20);
  \draw[Hedge] (n20) to[bend left=16] (n4);
  \draw[Sedge] (n20) to[bend left=16] (n12);
  \draw[Hedge] (n21) to[bend left=16] (n10);
  \draw[Sedge] (n21) to[bend left=16] (n23);
  \draw[Hedge] (n22) to[bend left=16] (n18);
  \draw[Sedge] (n22) to[bend left=16] (n9);
  \draw[Hedge] (n23) to[bend left=16] (n22);
  \draw[Sedge] (n23) to[bend left=16] (n14);
\end{tikzpicture}
  }
  \caption{Cayley graph of the quotient $\mathcal{C}_3/\mathcal{P}_3 \cong SL(2,\mathbb{Z}_3)$ generated by the traversal described in the text. Each of the $24$ nodes is one coset representative, labeled by the power-compressed word assigned to it when the breadth-first traversal first reached it. Orange arrows denote right multiplication by $H$ and purple arrows denote right multiplication by $S$. The $S$ arrows close into eight directed $3$-cycles and the $H$ arrows close into six directed $4$-cycles, reflecting the orders of the two generators, and every element is reachable from the identity $I$ within five generator applications.}
  \label{fig:cayley-c3modp3}
\end{figure}


\paragraph{Full expressions.}
The full catalog is the set of products $P \cdot C$, where $P = X^a Z^b$ with $a, b \in \mathbb{Z}_3$ ranges over the nine Pauli monomials and $C$ over the $24$ coset words. Identity factors are dropped, and the product of the two identities is the single word $I$. This yields exactly $216$ expression strings, such as $X^2Z^2S^2HSH$. Each string is a word in the symbols $X$, $Z$, $H$, $S$ and $I$, read as an operator product, so its rightmost symbol acts first and the Pauli prefix acts last. The strings contain between one and nine symbols, $5.17$ on average. Every Clifford of a benchmark sequence is drawn uniformly, with replacement, from these $216$ strings. For each of the $30$ sequences we draw one list of Cliffords as long as the largest depth of the run ($22$ in the production runs, $6$ in the high-$p$ extension runs of both noise models), and depth $m$ uses its first $m$ entries. LRB-D and terminal postselection (const-0) share these lists, physical RB draws its own, and every run draws new lists.

\paragraph{Physical realization.}
For physical RB, each symbol of a string is applied as its own primitive gate on the single qutrit, in the order obtained by reading the string from right to left. Under the circuit-level models of \autoref{sub:noise_models}, every primitive gate, each Pauli symbol included, is followed by its own DEP1 location, and the identity element occupies one idle moment. No gates are merged, so $X^2$ is two $X$ gates and $H^3$ is three $H$ gates.

\paragraph{Logical realization.}
For LRB-D, each symbol of a string is replaced by its logical circuit on the folded $[[5,1,2]]_3$ patch. The logical $\bar{H}$ and $\bar{S}$ are the constant-depth circuits of \autoref{fig: logical H} and \autoref{fig: logical S}. The logical Paulis $\bar{X} = X_0 \otimes X_3$ and $\bar{Z} = Z_0 \otimes Z_2$ are applied as two single-qutrit gates, $X$ on $q_0$ and $q_3$ or $Z$ on $q_0$ and $q_2$, which the scheduler places independently, so they need not share a moment. The logical identity is a single identity gate on $q_0$, which under circuit-level noise gives one moment of idle noise on every qutrit. As for physical RB, the substitution walks each string from right to left, so that the resulting circuit applies the factors in the correct temporal order. All physical gates of one logical Clifford form one scheduling segment of the noise model. Inside the segment, gates are packed as soon as possible into moments, so gates on disjoint qutrits share a moment. Consecutive gates on the same qutrit are never merged into a composite gate. Each primitive gate is a separate noisy operation, so the depth and the noise of a logical Clifford grow with the length of its word.

\paragraph{Inverses.}
The mirror inverse needs $C_i^{-1}$ for every sampled $C_i$. For a word $C = s_1 s_2 \cdots s_\ell$ the inverse is $C^{-1} = s_\ell^{-1} \cdots s_1^{-1}$, so in time order it applies $s_1^{-1}$ first and $s_\ell^{-1}$ last. The catalog stores this inverse word for each of its $216$ elements. It is written in $H^{-1}$, $S^{-1}$, $X^{-1}$ and $Z^{-1}$, has as many symbols as the forward word, and is not resynthesized as a shortest word. The mirror inverse looks up the stored inverse of each $C_i$ separately and realizes it like any other word, so every $C_i^{-1}$ is its own noisy Clifford. In LRB-D a logical inverse symbol is the time-reversed logical circuit with every gate inverted. For instance, $\bar{H}^{-1}$ applies $H$ to $q_1$ and $H^{-1}$ to $q_0, q_2, q_3, q_4$, followed by $\mathrm{SWAP}_{2,3}$.

\paragraph{Verification.}
The catalog and all circuits are built directly as \texttt{sdim} circuits, with $S$ realized by the \texttt{sdim} phase gate $\ket{j} \mapsto \omega^{j(j-1)/2} \ket{j}$. We check the catalog against the gate conventions of \texttt{sdim} by exporting single-qutrit circuits to Cirq and computing their unitaries. The $216$ strings give $216$ distinct unitaries modulo global phase, and they coincide with the closure of $\{H, S\}$ under multiplication. Every string followed by its stored inverse gives the identity up to phase. The same unitaries reproduce, in order, the $24$ coset words of the traversal and the $48$ edges of \autoref{fig:cayley-c3modp3}.
